\acte[La Triple Couronne d’or]{La Triple Couronne d’or du Sacré-Cœur de Jésus}
\usage{quot. : 1ʳᵉ couronne lun., mer., sam. ; 2ᵉ mar., ven. ; 3ᵉ jeu., dim.}
\rubrique{Première couronne, le lundi, le mercredi et le samedi. — Deuxième couronne, le mardi et le vendredi. — Troisième couronne, le jeudi et le dimanche. L’ordinaire ci-dessous est commun aux trois couronnes ; le propre de chacune suit.}

\intertitre{Invocations préparatoires}
\ouverture{Ô}{}{Jésus, mon doux Sauveur !} qui aimez tant les hommes, et qui cependant en êtes si peu aimé, permettez-moi, malgré mon indignité, de venir, en réparation des outrages et des ingratitudes dont sont payés vos bienfaits de chaque jour, vous offrir cette triple Couronne d’or. Je vous l’offre par le Cœur immaculé de votre très sainte Mère, pour la plus grande gloire du Père céleste et pour le règne de votre Cœur dans tous les cœurs.

Je vous l’offre en union à tous les bienheureux du paradis, à tous vos anges adorateurs de l’autel et à tous les associés de votre Garde d’honneur.

Je vous l’offre pour le monde chrétien, en particulier pour le salut de la France, et par elle pour la délivrance de notre sainte Mère l’Église.

Ô Cœur, flamme et lumière d’amour ! acceptez mon ardent désir de vous voir connu, aimé, consolé, glorifié partout ; à cette intention, bénissez mes paroles, mes actions, mes souffrances, acceptez mon cœur que je vous donne sans partage et sans retour, vous priant de le consumer du feu de votre charité. Ainsi soit-il.
\rubrique{On ajoute l’invocation propre de la couronne du jour.}

\intertitre{Introduction de la couronne}
\rubrique{Sur la croix :}
\priere\latin{Veni, Sancte Spiritus, reple tuorum corda fidelium, et tui amoris in eis ignem accende.}
\vr{\latin{Emitte Spiritum tuum, et creabuntur.}}{\latin{Et renovabis faciem terræ.}}
\priere\latin{Oremus. Deus, qui corda fidelium Sancti Spiritus illustratione docuisti, da nobis in eodem Spiritu recta sapere, et de ejus semper consolatione gaudere. Per Christum Dominum nostrum. Amen.}
\rubrique{Sur les cinq premiers grains, un Pater, suivi de l’invocation suivante, puis un Gloria :}
\priere\latin{Amatum sit ubique Sacratissimum Cor Jesu.}
\traduction{Aimé soit partout le Sacré-Cœur de Jésus.}

\intertitre{Corps de la couronne}
\rubrique{On annonce le mystère ; puis, sur le gros grain d’introduction :}
\vr{\latin{Jesu, mitis et humilis Corde,}}{\latin{fac cor meum sicut Cor tuum !}}
\traduction{Jésus, doux et humble de Cœur, rendez mon cœur semblable au vôtre.}
\indulg{\textit{Jesu, mitis et humilis Corde} : 500 j. ; plénière, cond. ord., si l’invocation est dite chaque jour pendant un mois. Ench.~1952, n\textsuperscript{o}~221.}
\rubrique{Sur la dizaine :}
\vr{\latin{Dulce Cor Jesu,}}{\latin{fac ut te magis ac magis diligam.}}
\traduction{Doux Cœur de Jésus, faites que je vous aime toujours de plus en plus.}
\rubrique{Sur le gros grain de conclusion :}
\vr{\latin{Dulce Cor Mariæ,}}{\latin{sis salus mea !}}
\traduction{Doux Cœur de Marie, soyez mon salut !}

\intertitre{Prières finales}
\rubrique{Consécration.}
\priere Très doux Cœur de Jésus, je vous consacre aujourd’hui et pour toujours mon esprit pour méditer vos perfections, mon cœur pour vous aimer, mon corps pour l’immoler à votre gloire.
\rubrique{Invocations.}
\priere Mon Jésus, miséricorde ! {\color{rubrique}\itshape (cinq fois)}
\priere Que partout et toujours soient connus, bénis, aimés, servis et glorifiés le très saint Cœur de Jésus et le Cœur immaculé de Marie !
\priere Saint Joseph, modèle et patron des amis du Sacré-Cœur de Jésus, priez pour nous.
\rubrique{Oblation.}
\priere\V Seigneur Jésus, daignez agréer l’oblation de nos cœurs, en réparation de tant et de si grandes injures qui sont faites à votre Sacré-Cœur, surtout dans le sacrement de votre amour, et faites-nous la grâce de ressentir davantage les douleurs de ce divin Cœur, d’imiter ses vertus et de mériter ses faveurs. Vous qui vivez et régnez dans les siècles des siècles.
\rubrique{Amende honorable.}
\priere\V Recevez, Père éternel, cette amende honorable que je vous fais, en union de celle que ce Sacré-Cœur vous fit sur le Calvaire, et de celle que Marie vous fit elle-même au pied de la Croix de son divin Fils.
\priere\R Ainsi soit-il.
\culdelampe

\acte[Première couronne]{Première couronne : le Sacré-Cœur de Jésus dans son incarnation}
\rubrique{Le lundi, le mercredi et le samedi.}
\intertitre{Invocation propre}
\priere Ô Cœur Sacré de Jésus, victime d’amour, je vous adore et je vous aime dans les mystères de votre incarnation.
\intertitre{Mystères}
\mystere{I.}{L’oblation du Sacré-Cœur de Jésus dans le sein virginal de Marie.}{offrir son cœur à Dieu dès le réveil, par le Sacré-Cœur de Jésus.}
\mystere{II.}{Le Sacré-Cœur de Jésus dans son enfance.}{la simplicité de cœur et l’obéissance à la sainte Église.}
\mystere{III.}{Le Sacré-Cœur de Jésus dans sa vie cachée à Nazareth.}{l’amour de sa condition et la fuite du monde.}
\mystere{IV.}{Le Sacré-Cœur de Jésus dans sa vie apostolique.}{union au Sacré-Cœur de Jésus dans son apostolat.}
\mystere{V.}{Le Sacré-Cœur de Jésus dévoué aux malades et aux pécheurs.}{amour de compassion aux malades, de dévouement aux pécheurs.}
\culdelampe

\acte[Deuxième couronne]{Deuxième couronne : le Sacré-Cœur de Jésus dans sa passion}
\rubrique{Le mardi et le vendredi.}
\intertitre{Invocation propre}
\priere Ô Cœur Sacré de Jésus, brisé de douleur à cause de nos péchés, immolez mon cœur à votre amour souffrant.
\intertitre{Mystères}
\mystere{I.}{Le Sacré-Cœur de Jésus dans son agonie.}{horreur et fuite du péché.}
\mystere{II.}{Le Sacré-Cœur de Jésus dans ses humiliations.}{douceur et patience dans les humiliations ; pardon généreux.}
\mystere{III.}{Le Sacré-Cœur de Jésus dans ses souffrances extérieures.}{la sanctification des souffrances par la résignation.}
\mystere{IV.}{Le Sacré-Cœur de Jésus abandonné du Ciel et de la terre.}{un abandon total aux bons plaisirs de Dieu.}
\mystere{V.}{Le Sacré-Cœur de Jésus percé par la lance.}{retraite de l’âme dans le Sacré-Cœur de Jésus.}
\culdelampe

\acte[Troisième couronne]{Troisième couronne : le Sacré-Cœur de Jésus dans l’Eucharistie}
\rubrique{Le jeudi et le dimanche.}
\intertitre{Invocation propre}
\priere Ô Cœur Sacré de Jésus, Cœur aimable et toujours aimant, vivant d’amour pour nous dans la sainte Eucharistie, que je ne vive plus que pour vous aimer !
\intertitre{Mystères}
\mystere{I.}{Vie d’amour.}{fidélité à visiter le Sacré-Cœur de Jésus au saint Sacrement et à assister, autant que possible, au saint sacrifice de la messe.}
\mystere{II.}{Vie de silence et de prière.}{le recueillement de l’âme dans le sein du Sacré-Cœur au milieu du monde.}
\mystere{III.}{Vie de sacrifice.}{le sacrifice pour Dieu et pour ses frères, à l’exemple du Sacré-Cœur de Jésus.}
\mystere{IV.}{Vie outragée par les méchants.}{faire chaque jour, au moins le vendredi, amende honorable au Sacré-Cœur de Jésus.}
\mystere{V.}{Vie d’actions de grâces.}{offrir sa reconnaissance au Cœur de Jésus par les saints Cœurs de Marie et de Joseph.}
\culdelampe